%%%%%%%%%%%%%%%%%%%%%%%%%%%%%%%%%%%%%%%%%%%%%%%%%%%%%%%%%%%%%%%%%%%
%% 
%% Arindam Chowdhury's resume
%%   - taken from Yisong Yue's webpage
%%   - based off work by Michael DeCorte 
%%
%%%%%%%%%%%%%%%%%%%%%%%%%%%%%%%%%%%%%%%%%%%%%%%%%%%%%%%%%%%%%%%%%%%



%%
%% The following code sets up the document formatting
%%

%this assumes that res_yy.sty is in some path
\documentstyle[hyperref, margin, line]{res_yy}

\hypersetup{backref,pdfpagemode=Full,colorlinks=true,backref}

\addtolength{\oddsidemargin}{-0.45in}
%\addtolength{\voffset}{0.10in}
\addtolength{\textwidth}{1.00in} %\addtolength{\textheight}{1.50in}

\renewcommand{\namefont}{\LARGE\emph}



%%
%% The following code defines some macros for terms which have raised font
%% (ie 4\fourth would result 4th with the 'th' raised (superscripted)
%%

\def\Cplusplus{{\rm C\raise.5ex\hbox{\small ++}}}
\def\CSharp{{\rm C\raise.5ex\hbox{\small \#}}}
% 'st' 'nd' 'rd' 'th' superscripts for numbers
\def\first{{\raise.5ex\hbox{\small st}}}
\def\second{{\raise.5ex\hbox{\small nd}}}
\def\third{{\raise.5ex\hbox{\small rd}}}
\def\fourth{{\raise.5ex\hbox{\small th}}}



%%
%% starting the actual document
%%

\begin{document}

%the name in big fonts at the top of resume
%this is left aligned
\name{Arindam Chowdhury}

%this is right aligned
\address{
website: \url{https://archo48.github.io/}\ \ \ \ \ 
email: ac131@rice.edu  
}

\begin{resume}



%%
%% This section of code is inelegant, but I'm too lazy to fix it
%%

\section{\textsc{Research Interest}}
My research interests lie primarily in signal procesing, machine learning and their applications in network science.

\section{\textsc{Education}}
\textbf{Ph.D in Data Science} \ \ \ \ \ \ \ \ \ \ \ \ \ \ \ \ \ \ \ \ \ \ \ \ \ \ \ \ \ \ \ \ \ \ \ \ \ \ \ \ \ \ \ \ \ \ \ \ \ \ \ \ \ \ \ \ \ \ \ \ \ \ \ \ \ \ \ \ \ \ \ \ \ \ \ \ \ \ \ \ \ \ 2019 - Present  \\
Rice University \ \ \ \ \ \ \ \ \ \ \ \ \ \ \ \ \ \ \ \ \ \ \ \ \ \ \ \ \ \ \ \ \ \ \ \ \ \ \ \ \ \ \ \ \ \\
\textit{CPI: 3.95/4.0}

\textbf{M.Tech in Signal Processing} \ \ \ \ \ \ \ \ \ \ \ \ \ \ \ \ \ \ \ \ \ \ \ \ \ \ \ \ \ \ \ \ \ \ \ \ \ \ \ \ \ \ \ \ \ \ \ \ \ \ \ \ \ \ \ \ \ \ \ \ \ \ \ \ \ \ 2015 - 2017 \\
Indian Institute of Technology (IIT) \ \ \ \ \ \ \ \ \ \ \ \ \ \ \ \ \ \ \ \ \ \ \ \ \ \ \ \ \ \ \ \ \ \ \ \ \ \ \ \ \ \ \ \ \\\
\textit{CPI: 9.32/10.0}

\textbf{B.Tech in Electronics and Communications} \ \ \ \ \ \ \ \ \ \ \ \ \ \ \ \ \ \ \ \ \ \ \ \ \ \ \ \ \ \ \ \ \ \ \ \ \ \ \ \ \ \ \ \ 2010 - 2014 \\ 
National Institute of Technology (NIT)\ \ \ \ \ \ \ \ \ \ \ \ \ \ \ \ \ \ \ \ \ \ \ \ \ \ \ \ \ \ \ \ \ \ \ \ \ \ \ \ \\ 
\textit{CPI: 8.70/10.0} 

%%
%% the meat of the resume starts now
%%

\begin{formatb}
  \employer{l}\title{r}\\
  \location{l}\dates{r}\\ 
  
  \body\\
\end{formatb}

\section{\textsc{Work Experience}}

\employer{\textbf{TCS Innovation Labs}}
\title{Researcher}
\location{Gurgaon, Haryana, India}
\dates{2017 - 2019}
\begin{position}
Information extraction from scanned images of documents. 
\end{position}

%\employer{\textbf{TCS Innovation Labs}}
%\title{Research Intern}
%\location{Gurgaon, Haryana, India}
%\dates{Summer 2016}
%\begin{position}
%Developed CNN based face recognition \& verification algorithms.
%\end{position}


%%
%% We use the same formatting for projects as for work experience
%% Shown below is the formatting used previously
%%
%%  \begin{formatb}
%%    \employer{l}\title{r}\\
%%    \location{l}\dates{r}\\
%%    \body\\
%%  \end{formatb}
%%
%% 
%%  Note that \location is now being used for non-location information
%%


\begin{formatb}
  \employer{l}\dates{r}\\
  \body\\
\end{formatb}

\section{\textsc{Projects}}

\employer{\textbf{Optimal Power Allocation in Wireless Networks}}
\dates{Spring 2020 - Present}
\begin{position}
Combining a graph convolutional network (GCN) based learning system with an iterative algorithm (WMMSE) that incorporates domain structure to develop a fast, efficient and effective framework for power allocation in wireless networks.
\end{position}

\employer{\textbf{Efficient Training of Graph Convolutional Networks}}
\dates{Fall 2020 - Present}
\begin{position}
Developing a distributed framework for efficient training of graph convolutional networks (GCN) through partitioning of the hidden layers with a provision for graph clustering to facilitate mini-batch SGD in case of large graphs.
\end{position}


%\employer{\textbf{Molecular Graph Generation}}
%\dates{Summer 2020 - Present}
%\begin{position}
% 
%\end{position}

\employer{\textbf{Document Image Analysis}}
\dates{2017 - 2019}
\begin{position}
Automatic generation of textual summary of graphical data through reasoning-based question-answering. End-to-end neural model for offline recognition of handwritten text. Automatic digitization of industrial inspection sheets.
 
\end{position}


%\newpage
\section{\textsc{Publications}}

\textit{Unfolding wmmse using graph neural networks for efficient power allocation}\\
Under review at IEEE Transactions on Wireless Communications (\textbf{TWC})\\
\textbf{A. Chowdhury}, G. Verma, C. Rao, A. Swami and S. Segarra\\
\url{https://arxiv.org/pdf/2009.10812}

\textit{GIST: Distributed Training for Large-Scale Graph Convolutional Networks}\\
Under review at International Conference on Machine Learning (\textbf{ICML}) 2021\\
C. R. Wolfe, J. Yang, \textbf{A. Chowdhury}, C. Dun, A. Bayer, S. Segarra and A. Kyrillidis

\textit{Efficient power allocation using graph neural networks and deep algorithm unfolding}\\
Accepted at IEEE International Conference on Acoustics, Speech and Signal Processing (\textbf{ICASSP}) 2021\\ 
\textbf{A. Chowdhury}, G. Verma, C. Rao, A. Swami and S. Segarra\\  \url{https://arxiv.org/pdf/2012.02250},

\textit{ChartNet: Visual Reasoning over Statistical Charts using MAC-Networks}\\
Published at International Joint Conference on Neural Networks (\textbf{IJCNN}) 2019\\
M. Sharma, S. Gupta, \textbf{A. Chowdhury} and L. Vig\\
\url{https://ieeexplore.ieee.org/document/8852427}

\textit{An Efficient End-to-End Neural Model for Handwritten Text Recognition}\\
Published at British Machine Vision Conference (\textbf{BMVC}) 2018\\
\textbf{A. Chowdhury} and L. Vig\\
\url{http://bmvc2018.org/contents/papers/0606.pdf}

\textit{Reading Industrial Inspection Sheets by Inferring Visual Relations}\\
Published at IWRR Asian Conference on Computer Vision (\textbf{ACCV}) 2018\\
R. Rahul, \textbf{A. Chowdhury}, Animesh, S. Mittal and L. Vig\\
\url{https://arxiv.org/pdf/1812.07104.pdf}

%\textit{Deep Reader: End-to-end framework for relevant text extraction from real image documents with Natural Language Interface}\\
%Accepted at IWRR Asian Conference on Computer Vision (\textbf{ACCV}) 2018\\
%Vishwa D., R. Rahul, G. Sehgal, Swati, \textbf{A. Chowdhury}, M. Sharma, L. Vig, G. Shroff and A. Srinivasan\\
%\url{https://arxiv.org/pdf/1812.04377.pdf}

%
%\section{\textsc{Patents}}
%
%\textbf{A. Chowdhury} and L. Vig.\\ \textit{Systems and Methods for End-to-End Handwritten Text Recognition using Neural Networks}. Indian Patent Application 201821026934, filed July 2018. Patent pending.
%
%R. Rahul, \textbf{A. Chowdhury}, Animesh, S. Mittal and L. Vig.\\ \textit{Digitization of Industrial Inspection Sheets by
%Inferring Visual Relations}.\\ Indian Patent Application 201821044939, filed November 2018. Patent pending.
%
%Vishw, Rohit, Gunjan, \textbf{A. Chowdhury},  Swati, Monika, L. Vig, G. Shroff, and A. Srinivasan. \textit{Method and System for Information Extraction from Document Images using Conversational Interface and Database Querying}.\\ Indian Patent Application 201821045427, filed November 2018. Patent pending.

%\newpage
\section{\textsc{Course Work}}
\begin{tabular}{llll}
Statistical Signal Processing & \ \ &  Network Science and Analytics\\
Optimization: Algorithms, Complexity \& Approximation  & \ \ & Random Processes \\
Linear Algebra \& Optimization   & \ \ & Multi-Agent Dynamic Systems \\  
Pattern Recognition \& Machine Learning   & \ \ & Information Theory \\
Signal Processing Algorithms \& Architectures   & \ \ & Computer Vision \\
\end{tabular}

\section{\textsc{Technical Skills}}

\emph{Programming Languages}: C, \Cplusplus, Python\\
\emph{Numerical Computing Tool}:  MATLAB, Octave\\
\emph{Libraries}: dlib, NetworkX, Tensorflow, Keras, PyTorch\\
\emph{Cloud Computing}: Amazon EC2, Google Colab\\
\emph{Operating Systems}: Ubuntu 18.04 LTS, Windows
     
\section{\textsc{Academic Activities}}

\begin{formatb}
  \employer{l}\dates{r}\\
  \body\\
\end{formatb}

%\employer{\textbf{ACM CODS-COMAD}}
%\dates{Spring 2018}
%\begin{position}
%Participated in ACM India Joint International Conference on Data Science and Management of Data, 2018 edition held in Goa.
%\end{position}
%\employer{\textbf{ACCV}}
%\dates{December 2018}
%\begin{position}
%Presented orally in IWRR Asian Conference on Computer Vision, 2018 edition held in Perth, Australia. 
%\end{position}
%
%\employer{\textbf{BMVC}}
%\dates{October 2018}
%\begin{position}
%Presented poster in British Machine Vision Conference, 2018 edition held in Newcastle, England. 
%\end{position}
\employer{\textbf{Conference Reviewer}}
\dates{Spring 2020 - Present}
\begin{position}
IEEE ISIT 2021, IEEE ICASSP 2021, IEEE ASILOMAR 2020, IEEE VTC 2020  
\end{position}

\employer{\textbf{Rice University Teaching Assistant}}
\dates{Spring 2021}
\begin{position}
Grader for a course on Signals, Systems, and Transforms
\end{position}

\employer{\textbf{Rice University Teaching Assistant}}
\dates{Fall 2020}
\begin{position}
Grader for a course on Network Science and Analytics
\end{position}

\employer{\textbf{Rice University Teaching Assistant}}
\dates{Fall 2019}
\begin{position}
Grader for a course on Data and Dynamical Systems
\end{position}


%\employer{\textbf{IIT Guwahati Teaching Assistant}}
%\dates{Fall 2016}
%\begin{position}
%Prepared \& graded assignments for a course on Digital Signal Processing.
%\end{position}

%\employer{\textbf{IIT Guwahati Teaching Assistant}}
%\dates{Spring 2016 \& Spring 20117}
%\begin{position}
%Served as the instructor for Basic Electronics Laboratory course twice.
%\end{position}

%\employer{\textbf{ICVGIP}}
%\dates{December 2016}
%\begin{position}
%Organized, volunteered and participated in Indian Conference on Computer Vision, Graphics \& Image Processing, 2016 edition held at IIT Guwahati. 
%\end{position}
%
%\employer{\textbf{Deep Learning Summer School}}
%\dates{Summer 2016}
%\begin{position}
%Participated in a 6-day Summer School On Deep Learning for Computer Vision organized by IIIT Hyderabad
%\end{position}

% \section{Referees}
% \begin{itemize}
% \item[1]\href{https://in.linkedin.com/in/lovekesh-vig-5ba6ba2}{\textbf{\underline{Dr. Lovekesh Vig}}} \\ \\
% Senior Scientist, \\ 
% TCS Innovation Labs, Gurgaon, India. \\
% %Phone: +91-361-2582523 (O) \\
% Email: lovekesh.vig@tcs.com
%  \\
% \item[2]\href{https://in.linkedin.com/in/gautam-shroff-066901}{\textbf{\underline{Dr. Gautam Shroff}}} \\ \\
% Chief Scientist \& Head of Research,\\ 
% TCS Innovation Labs, Gurgaon, India. \\
% %Phone: +91-361-2582523 (O) \\
% Email: gautam.shroff@tcs.com
% \\ 
%% \item[4] \href{https://www.iitg.ernet.in/eee/pguha.html}{\textbf{\underline{Dr. Prithwijit Guha}}} \\ \\
%% Assistant Professor, EEE Department, \\
%% IIT Guwahati, Guwahati-781039, \\
%% Assam, India. \\
%% Phone: +91-361-2583452 (O) \\
%% Email: pguha@iitg.ernet.in
%%  \\ \\
% \item[3]\href{https://www.iitg.ernet.in/eee/mkb.html}{\textbf{\underline{Dr. Manas Kamal Bhuyan}}} \\ \\
% Professor, EEE Department, \\
% IIT Guwahati, Assam, India \\
% %Phone: +91-361-2582523 (O) \\
% Email: mkb@iitg.ernet.in
% \end{itemize}

\end{resume}
\end{document}
